\subsection{Generating the JAX model}

Once the session has passed the checks described above, generate the executable
JAX model by running:

\begin{verbatim}
pfit jax sessions/<session-name>
\end{verbatim}

This command converts the readable NumPy model in
\texttt{generated/user\_model.py} into the differentiable JAX implementation
used for parameter estimation. The translation is performed with the configured
local language model, but it is not treated as a single unrestricted code
generation step. The workflow asks for smaller fragments: helper functions,
right-hand-side expressions, the loss body when needed, and the writeout body
when needed. Source-derived or deterministically generated fragments are
protected from unrelated repair attempts.

The resulting script is written to:

\begin{verbatim}
sessions/<session-name>/generated/generated_script.py
\end{verbatim}

For the running example, the generated script contains the differential
equations for drug amount and response, interpolation of
\texttt{dose\_rate}, calculation of \texttt{signal\_model}, the loss function
defined in the input statement, and the columns written to the fitted-output
files.

A simplified excerpt of the transformed JAX system for this example is:

\begin{tcblisting}{
  listing only,
  breakable,
  colback=gray!6,
  colframe=gray!35,
  boxrule=0.4pt,
  arc=1mm,
  left=2mm,
  right=2mm,
  top=1mm,
  bottom=1mm
}
@jax.jit
def user_defined_system(t, y, other_args):
    constants = other_args["constants"]
    trainable_variables = other_args["trainable_variables"]
    dataset = constants["dataset"]
    t_eval = constants["t_eval"]

    k_elim = unscale_value(trainable_variables[0], 0.01, 10.0, True)
    k_on = unscale_value(trainable_variables[1], 0.001, 10.0, True)
    k_off = unscale_value(trainable_variables[2], 0.001, 10.0, True)
    baseline = unscale_value(trainable_variables[3], 50.0, 150.0, False)
    gain = unscale_value(trainable_variables[4], 1.0, 200.0, False)

    D = y[0]
    R = y[1]
    dose_rate = jnp.interp(t, t_eval, dataset[:, 0])

    dDdt = dose_rate - k_elim * D
    dRdt = k_on * D * (1.0 - R) - k_off * R
    return jnp.array([dDdt, dRdt])

@jax.jit
def _compute_loss_value(constants, trainable_variables,
                        solution_time, solution):
    dataset = constants["dataset"]
    baseline = unscale_value(trainable_variables[3], 50.0, 150.0, False)
    gain = unscale_value(trainable_variables[4], 1.0, 200.0, False)

    R = solution[:, 1]
    signal_model = baseline + gain * R
    signal = dataset[:, 1]
    mask = jnp.isfinite(signal)
    scale = jnp.max(jnp.abs(signal[mask]))
    residual = (signal_model[mask] - signal[mask]) / scale
    return jnp.sqrt(jnp.mean(residual**2))
\end{tcblisting}

This excerpt omits the surrounding integration scaffold, error handling,
writeout routine, and multi-experiment aggregation, but it shows the main
transformation: parameters are mapped from optimizer coordinates to physical
values, the forcing input is interpolated from the non-time dataset columns,
the right-hand side is written with \texttt{jax.numpy}, and the stated
normalized RMSE is evaluated on finite measurements.

The generated script is not accepted simply because it is valid Python. The
framework first checks that it provides the required runtime interface for
integration, loss evaluation, and output writing. It then imports the script,
runs the model for each dataset, verifies that trajectories and losses are
finite, and evaluates the generated code on sampled parameter points. These
sampled points are chosen in the normalized search coordinates used by the
optimizers, including logarithmic coordinates for parameters marked as
log-scaled.

The JAX implementation is compared with the readable NumPy source model at
sampled states and parameter values. The comparison checks differential-equation
values, observable predictions, output columns, and losses where applicable.
These tests are intended to detect accidental changes introduced during
translation, such as altered parameter names, incorrect column indices, missing
observable transformations, or a loss that no longer matches the source model.
They provide useful safeguards, but they do not prove equivalence for every
possible state and parameter value.

\paragraph{Numerical readiness checks.}
The JAX stage also performs numerical checks needed before fitting. For fresh
generated configurations, the framework may select an integrator from the data,
estimate an initial \texttt{max\_steps} budget, test solver coverage across
sampled parameter points, calibrate state-scaled tolerances, and compare
predictions and losses against integrations run with tighter tolerances. These
checks are deterministic and are recorded in JSON reports under
\texttt{generated/}.

The principal reports include:

\begin{verbatim}
solver_selection.json
solver_coverage.json
solver_diagnostics.json
solver_recovery.json
tolerance_calibration.json
solver_accuracy.json
translation_fidelity.json
\end{verbatim}

If too few sampled parameter points integrate successfully because the solver
reaches its step limit, the workflow may increase \texttt{max\_steps} up to the
configured recovery ceiling and rerun the same validation. This recovery changes
only the numerical step budget; it does not ask the model to rewrite the
scientific equations, loss, or observables. If recovery succeeds, the accepted
step budget is written back to the YAML configuration. If recovery is exhausted
or times out, the session is not marked ready.

Tolerance calibration is similarly guarded. Proposed state-scaled tolerances or
population-stage relaxations are accepted only when they preserve finite losses
and the meaningful ordering of sampled candidates. Solver-accuracy checks then
compare selected feasible probes with tighter reference solves. When the
accuracy check rejects a proposed relaxation, the workflow falls back to the
validated stricter profile rather than accepting an inaccurate setting.

\paragraph{Repair and source freshness.}
If a translation comparison or interface check fails, the framework may ask the
local model to repair only the fragment responsible for the failure. The number
of repair attempts is limited, and repairs are field-scoped so that a loss
failure does not permit arbitrary changes to equations, parameter definitions,
or protected source-derived bodies. A session is not marked ready if the
generated model still fails its checks.

The framework records which versions of \texttt{inputs/user\_input.yaml} and
\texttt{generated/user\_model.py} were used to create the JAX script. If either
file is edited afterwards, the generated script is treated as out of date and
\texttt{pfit jax} must be run again before fitting.

Before fitting, the user should inspect \texttt{generated\_script.py} and the
generated reports and confirm that the equations, parameter names, forcing
interpolation, observable definitions, loss function, output columns, integrator,
tolerances, and step budget match the intended model.

Finally, verify that the complete session is ready to run:

\begin{verbatim}
pfit check sessions/<session-name> --ready
\end{verbatim}
